\section{Background}
\label{appendix:background}
\textbf{Partially Observable Markov Decision Process.}
% A Markov Decision Process (MDP) is described by a tuple $(S, A, P, R)$, with $s_t \in S$ is a state, $a_t \in A$ is an action, $p(s' | s_t = s, a_t = a)$ is a probability of transition from state $s_t$ to state $s_{t+1}$ after taking an action $a_t$, $R(s, a)$ is a reward for taking an action $a$ at the state $s$. At each timestep $t$ an agent $\pi$ observes a state $s_t$, selects an action $a_t \sim \pi(\cdot|s_t)$ and observes the next state $s_{t+1} \sim P(\cdot | s_t, a_t)$ and a reward $R(s_t, a_t)$. In Partially Observable MDPs instead of a state $s$, and agent receives an observation $o_t$ which carries a partial information about the real state of the MDP. In the context of our work, the observation might miss the goal information, and thus this information must be inferred from the agent's memory.
A Markov Decision Process (MDP) is defined by a tuple $(S, A, P, R)$, where $s_t \in S$ denotes a state, $a_t \in A$ an action, $p(s' | s_t = s, a_t = a)$ the transition probability from state $s_t$ to $s_{t+1}$ after taking action $a_t$, and $R(s, a)$ the reward for action $a$ in state $s$ \citep{sutton2018reinforcement}. An agent $\pi$ observes the state $s_t$, selects action $a_t \sim \pi(\cdot|s_t)$, and receives the subsequent state $s_{t+1} \sim P(\cdot | s_t, a_t)$ and reward $R(s_t, a_t)$. In POMDPs, the agent receives an observation $o_t$ instead of the full state $s$, which contains partial information about the MDP's real state. In the context of our work, $o_t$ may lack goal information, requiring inference from the agent's memory.

\textbf{Multi-Armed Bandits.}
% \textbf{Bernoulli.} A Multi-Armed Bandit (MAB) is an environment with $N$ arms $a_i \in A$. Each arm is associated with a mean value $\mu_i$ and pulling each arm results in a reward $r_i \in R$ generated from $\text{Bernoulli}(\mu)$. The agent's goal is to efficiently find the best arm with the highest $\mu$. To evaluate the model's performance, we use the standard bandit metric - regret. A regret at timestep t is defined as $\sum_t(\mu_{a^*} - \mu_{\hat{a}_t})$.
A Bernoulli multi-armed bandit (MAB) environment consists of $N$ arms $a_i \in A$, each associated with a mean $\mu_i$ \citep{sutton2018reinforcement}. Pulling an arm $a_i$ yields a reward $r_i \sim \text{Bernoulli}(\mu_i)$. The agent's objective is to identify the arm with the highest $\mu_i$. Performance is measured using regret, calculated as $\sum_t(\mu_{a^*} - \mu_{\hat{a}_t})$. Unlike MDPs, MABs lack states. 

In a Contextual MAB, each arm $a$ has a feature vector $x_a \in R^d$ \citep{sutton2018reinforcement}. At each step $t$, the agent observes a context state $s_t$. Selecting action $a_i$ results in a reward from a normal distribution with mean $\mu_t = \langle s_t, a_i \rangle$ and standard deviation $\sigma$. Here, unlike in MDPs, the actions influence immediate rewards but not future states.

% Compared to an MDP, an MAB does not have states.

% \textbf{Contextual.} A Contextual MAB operates as follows. Each arm $a$ has a feature vector $x_a \in R^d$. At each step $t$ an agent observes a state (\textit{context}) $s_t$. Choosing an action $a_i$ will return a reward generated by a normal distribution with a mean $\mu = <s_t, a_i>$ and a standard deviation $\sigma$. 

% Compared to an MDP, in a Contextual MAB current action does not affect future states, only immediate rewards.

\textbf{In-Context Learning.} In-Context Learning describes the capability of a model to infer its task from the context it is given. For instance, the GPT-3 model \citep{brown2020language} can be prompted in natural language to perform a variety of functions such as text classification, summarization, and translation, despite not being explicitly trained for these specific tasks. One of the possible prompts, which is used in our paper in some form, is a list of example pairs $(x_i, y_i)$ ending with a query $x_q$ for which the model is expected to generate a corresponding prediction $\hat{y}_q$.

\textbf{Contrastive Learning.} Contrastive learning focuses on creating representations where “similar” examples are close together in the feature space, while “dissimilar” ones are far apart \citep{weinberger2009distance, schroff2015facenet}. This concept may be encapsulated in a triplet loss formula, where the similarity between an anchor and a positive example is maximized, and that between an anchor and a negative example is minimized, expressed as $L = sim(x, x^+) - sim(x,x^-)$. \citet{oord2018representation} has developed a variant of contrastive loss called InfoNCE. This loss was lately widely adopted for representation learning \citep{jaiswal2020survey}.
